%% 摘要

% 中文摘要
\begin{abstract}
    近年来，随着航天强国建设持续推进，我国深空探测任务不断向月球、火星和更远行星系统拓展，获取了海量的深空观测数据，这些数据在行星科学研究、地外天体环境理解、空间资源开发利用及任务规划等方面发挥着越来越重要的作用。然而，传统的二维观测数据难以直观、连续地展现地外天体复杂的空间几何与光影特征，无法满足深空任务对高精度空间认知和沉浸式智能感知的前沿需求。为此，深空场景的三维重建与新视角合成成为深空探测数据智能分析和三维可视化的核心环节。该技术旨在从有限多视角图像中恢复可渲染的三维场景表示，并在未观测视角下生成连续、稳定的渲染图像。然而，与地球常规环境相比，深空场景通常存在观测视角有限、成像距离远、远近深度跨度大、地表纹理弱以及光照条件不稳定等问题。这些因素削弱了多视角图像中的几何约束和光度一致性，使现有面向常规地面场景的三维重建与新视角合成方法难以直接迁移应用。

    针对上述问题，本文提出一种面向深空场景的Deep-Space Scene Gaussian Splatting（DSS-GS）方法。该方法以3D Gaussian Splatting为基础，构建显式三维高斯场景表示，并通过可微渲染完成场景优化与新视角合成。为适应深空场景中的远区域弱监督、高斯形态不稳定以及跨视角光度不一致问题，本文在原始3DGS框架中引入远区域增强策略、高斯形态稳定约束和可选的参考引导外观一致性模块，使模型能够在保持整体渲染质量的同时，更充分地表达远处地形区域和弱纹理结构，并缓解光照扰动导致的跨视角色彩偏差。
        
    本文的主要工作和成果如下：

    （1）将一种基于深度先验的远区域增强策略引入3DGS。Far-region Adaptive Mask Enhancement（FAME）利用辅助单目深度构建自适应远区域掩膜，从而增强远距离地形区域在图像域中的监督信号，缓解远区域在优化过程中容易监督不足的问题。

    （2）提出一种高斯形态稳定化模块，用于约束显式三维高斯基元。Rank-Volume Stabilizer（RVS）结合有效秩正则和体积正则，对高斯形态分布和空间覆盖范围进行约束，从而提升高斯基元在有限视角和弱纹理条件下的表示稳定性。

    （3）设计了一种基于参考引导的外观一致性可选模块，作为跨视角光度不一致条件下的辅助处理路径，用于建模视角间亮度和色彩差异，缓解光照扰动对新视角合成结果的影响。

    （4）在M3arsSynth火星地表序列上的实验结果表明，DSS-GS取得了0.961的SSIM、38.49 dB的PSNR和0.078的LPIPS。与3DGS基线相比，DSS-GS的SSIM提升0.030，PSNR提升2.32 dB，LPIPS降低0.053，说明本文方法能够有效提升深空地表场景的新视角合成质量。

    \keywordslist{三维重建, 新视角合成, 3D Gaussian Splatting, 深空场景}
\end{abstract}

% 英文摘要
\begin{engabstract}
    In recent years, with the continuous advancement of space exploration, deep-space missions have expanded from the Moon and Mars toward more distant planetary systems, producing massive volumes of deep-space observational data. These data are playing an increasingly important role in planetary science, extraterrestrial environment understanding, space resource exploitation, and mission planning. However, conventional two-dimensional observations cannot intuitively and continuously characterize the complex spatial geometry, illumination, and shading properties of extraterrestrial bodies, and thus cannot fully meet the emerging demands of deep-space missions for high-precision spatial cognition and immersive intelligent perception. Consequently, 3D reconstruction and novel view synthesis for deep-space scenes have become key components of intelligent analysis and 3D visualization of deep-space exploration data. The goal is to recover a renderable 3D scene representation from limited multi-view images and generate continuous and stable renderings from unobserved viewpoints. Compared with ordinary terrestrial environments, deep-space scenes often suffer from limited viewing coverage, long-range imaging, large near-far depth variations, weak surface textures, and unstable illumination. These factors weaken geometric constraints and photometric consistency in multi-view imagery, making it difficult to directly transfer existing 3D reconstruction and novel view synthesis methods designed for conventional ground-level scenes.

    To address these challenges, this thesis proposes Deep-Space Scene Gaussian Splatting (DSS-GS). Built upon 3D Gaussian Splatting (3DGS), DSS-GS constructs an explicit 3D Gaussian scene representation and optimizes the scene through differentiable rendering for novel view synthesis. To adapt 3DGS to weak supervision in distant regions, unstable Gaussian morphology, and cross-view photometric inconsistency in deep-space scenes, this thesis introduces a far-region enhancement strategy, a Gaussian morphology stabilization constraint, and an optional reference-guided appearance consistency module into the original 3DGS framework. These designs enable the model to more sufficiently represent distant terrain regions and weakly textured structures while preserving the overall rendering quality, and they alleviate cross-view color deviations caused by illumination perturbations.
        
    The main contributions and results of this thesis are summarized as follows:

    (1) A depth-prior-based far-region enhancement strategy is introduced into 3DGS. Far-region Adaptive Mask Enhancement (FAME) uses auxiliary monocular depth to construct an adaptive far-region mask, thereby strengthening image-domain supervision for distant terrain regions and alleviating the under-supervision problem that commonly occurs during optimization.

    (2) A Gaussian morphology stabilization module is proposed to regularize explicit 3D Gaussian primitives. Rank-Volume Stabilizer (RVS) combines effective-rank regularization with volume regularization to constrain the shape distribution and spatial coverage of Gaussian primitives, thereby improving their representation stability under limited-view and weak-texture conditions.

    (3) An optional reference-guided appearance consistency module is designed as an auxiliary processing path for cross-view photometric inconsistency. This module models brightness and color differences across viewpoints and mitigates the influence of illumination perturbations on novel view synthesis.

    (4) Experiments on Martian surface sequences from M3arsSynth show that DSS-GS achieves an SSIM of 0.961, a PSNR of 38.49 dB, and an LPIPS of 0.078. Compared with the 3DGS baseline, DSS-GS improves SSIM by 0.030 and PSNR by 2.32 dB while reducing LPIPS by 0.053, demonstrating that the proposed method effectively improves the quality of novel view synthesis for deep-space surface scenes.

    \engkeywordslist{3D reconstruction, novel view synthesis, 3D Gaussian Splatting, deep-space scenes}
\end{engabstract}
